% Part: computability
% Chapter: computability-theory
% Section: total

\documentclass[../../../include/open-logic-section]{subfiles}

\begin{document}

\olfileid[fr]{cmp}{thy}{tot}
\olsection{La totalité est indécidable}

Voyons un autre emploi du théorème $s$-$m$-$n$
pour établir qu'un ensemble n'est pas calculable.
Notons $\fn{Tot}$ l'ensemble des indices des
fonctions calculables totales, c'est-à-dire
\[
\fn{Tot} = \Setabs{x}{\text{pour tout $y$, $\cfind{x}(y)\fdefined$}}.
\]

\begin{prop}
\ollabel{prop:total}
$\fn{Tot}$ n'est pas calculable.
\end{prop}

\begin{proof}
Il suffit de réduire $K$ à $\fn{Tot}$.
Définissons $h(x,y)$ par
\[
h(x,y) \simeq
\begin{cases}
0 & \text{si $x \in K$} \\
\fundefined & \text{sinon}
\end{cases}
\]
La valeur de $h(x,y)$ ne dépend pas de~$y$.
Cette fonction est calculable partielle : sur
l'entrée $(x,y)$, on simule $\cfind{x}$ sur
l'entrée~$x$; si ce calcul s'arrête, on renvoie~$0$,
et sinon le calcul ne s'arrête pas. Ainsi
$h(x,y)\simeq\Zero(\umin{s}{T(x,x,s)})$,
où $\Zero$ est la fonction constante égale à zéro. La composition de fonctions partielles exige que la recherche intérieure termine : si aucun témoin $s$ n'existe, cette expression reste indéfinie, même si $\Zero$ ignore la valeur de son argument.

Le théorème $s$-$m$-$n$ fournit une fonction
récursive primitive~$k(x)$ telle que, pour tous
$x$ et~$y$,
\[
\cfind{k(x)}(y) \simeq
\begin{cases}
0 & \text{si $x \in K$} \\
\fundefined & \text{sinon}
\end{cases}
\]
La fonction $\cfind{k(x)}$ est donc totale si
$x\in K$, et n'est définie nulle part sinon.
Par conséquent, $k$ réduit $K$ à~$\fn{Tot}$.
\end{proof}

\begin{digress}
En fait, $\fn{Tot}$ n'est même pas effectivement
énumérable : sa complexité se situe plus haut
dans la « hiérarchie arithmétique ». Nous
n'aurons pas besoin de ce renforcement ici.
\end{digress}

\end{document}
